\title{Επισκόπηση Λειτουργίας Δικτύων Διανομής Ηλεκτρικής Ενέργειας}

\section{Εισαγωγή}
\label{kef-04-theoria}

Σκοπός των δικτύων διανομής ηλεκτρικής ενέργειας (ΔΔΗΕ) είναι η ασφαλής και συνεχής
τροφοδότηση των τελικών χρηστών. Για δεκαετίες τα δίκτυα αυτά σχεδιάζονταν και λειτουργούσαν
ως \emph{παθητικά}, με την παραδοχή μονόδρομης ροής ισχύος από τον υποσταθμό ΥΤ/ΜΤ
προς τα φορτία. Επιπλέον, η τάση ρυθμιζόταν μία φορά, κατά τη μελέτη, ενώ η εποπτεία
περιοριζόταν στον υποσταθμό. Ωστόσο, η διείσδυση των διανεμημένων ενεργειακών πόρων (ΔΕΠ),
όπως τα φωτοβολταϊκά και οι σταθμοί αποθήκευσης, σε συνδυασμό με τον εξηλεκτρισμό των
μεταφορών και της θέρμανσης, έχει μετατρέψει τα δίκτυα διανομής σε \emph{ενεργά} δίκτυα, στα
οποία η ροή ισχύος αντιστρέφεται κατά διαστήματα και η τάση μεταβάλλεται ανάλογα με τις συνθήκες φόρτισης του δικτύου, με αποτέλεσμα ο διαχειριστής να χρειάζεται περισσότερα μέσα εποπτείας και
ελέγχου, συχνά σε πραγματικό χρόνο.

Στην Ελλάδα το Ελληνικό Δίκτυο Διανομής Ηλεκτρικής Ενέργειας (ΕΔΔΗΕ) λειτουργεί ο
Διαχειριστής Ελληνικού Δικτύου Διανομής Ηλεκτρικής Ενέργειας (ΔΕΔΔΗΕ), ο οποίος είναι
υπεύθυνος για την ανάπτυξη, τη λειτουργία και τη συντήρηση του δικτύου, τη σύνδεση των χρηστών
και τις μετρήσεις, και στα μη διασυνδεδεμένα νησιά ασκεί και καθήκοντα διαχειριστή συστήματος.
Το κεφάλαιο περιγράφει συνοπτικά τη δομή των δικτύων διανομής και τις απαιτήσεις λειτουργίας
τους και εστιάζει στην εποπτεία και τον έλεγχό τους: στους έξυπνους μετρητές και τους
ψηφιακούς υποσταθμούς, στα συστήματα SCADA και ADMS, στην εκτίμηση κατάστασης και στην
αυτόματη αποκατάσταση της τροφοδότησης. Η ρύθμιση της τάσης, η διαχείριση των συμφορήσεων και ο
συντονισμός με τον ΔΣΜΗΕ εξετάζονται στα
Κεφάλαια 5, 6 και
7.

\section{Δομή των δικτύων διανομής}
\label{sec:dn-structure}

\subsection{Επίπεδα τάσης και υποσταθμοί}
\label{sec:dn-levels}

Το δίκτυο διανομής αρχίζει στους \emph{υποσταθμούς ΥΤ/ΜΤ} (πρωτεύοντες υποσταθμοί), όπου η
υψηλή τάση των 150 kV του συστήματος μεταφοράς υποβιβάζεται στη μέση τάση, και περιλαμβάνει το
δίκτυο μέσης τάσης (ΜΤ), τους \emph{υποσταθμούς ΜΤ/ΧΤ} (δευτερεύοντες υποσταθμοί ή υποσταθμοί
διανομής) και το δίκτυο χαμηλής τάσης (ΧΤ) έως τον μετρητή κάθε καταναλωτή
(Σχήμα~\ref{fig:dn-structure}). Ο υποσταθμός ΥΤ/ΜΤ τροφοδοτείται από το σύστημα μεταφοράς και,
μέσω μετασχηματιστών με ΣΑΤΥΦ, τροφοδοτεί τους ζυγούς ΜΤ, από τους οποίους αναχωρούν οι γραμμές
ΜΤ. Το δίκτυο είναι ως επί το πλείστον εναέριο στις αγροτικές περιοχές και συνήθως υπόγειο στις
αστικές. Όσον αφορά την τοπολογία, το δίκτυο είναι ακτινικό, εκτός από τις αστικές περιοχές,
όπου σχεδιάζεται με βρόχους ώστε να υπάρχει και εναλλακτική διαδρομή τροφοδότησης· οι βρόχοι
όμως λειτουργούν ακτινικά, μέσω κανονικά ανοικτών διακοπτών (ΚΑ). Από κάθε υποσταθμό ΜΤ/ΧΤ
αναχωρούν γραμμές ΧΤ προς τους καταναλωτές, καθένας από τους οποίους συνδέεται με δική του
μονοφασική ή τριφασική παροχή. Στο ΕΔΔΗΕ η ΜΤ είναι κατά κανόνα 20 kV και η ΧΤ 400 V. Ο
ΔΕΔΔΗΕ κατέχει επίσης μικρό μήκος γραμμών ΥΤ στην Αττική και στα νησιά.

\begin{figure}[!htbp]
\centering
\includegraphics[width=\linewidth]{figures/dn-structure.svg}
\caption{Δομή ενός δικτύου διανομής.}
\label{fig:dn-structure}
\end{figure}

Ο υποσταθμός ΥΤ/ΜΤ περιλαμβάνει συνήθως μετασχηματιστές 150/21 kV ονομαστικής ισχύος 40/50 MVA
ή, σε περιοχές μικρότερου φορτίου, 20/25 MVA. Το δευτερεύον τους έχει ονομαστική τάση 21 kV,
δηλαδή 5\% πάνω από την ονομαστική τάση του δικτύου, ώστε στην ονομαστική λήψη η τάση των ζυγών
ΜΤ να αντισταθμίζει μέρος της πτώσης τάσης των γραμμών. Το ΣΑΤΥΦ βρίσκεται στο τύλιγμα της ΥΤ
και έχει 17 θέσεις με βήμα 1,875 kV, δηλαδή 1,25\% των 150 kV, σε εύρος από +7,5\% έως −12,5\%.
Η ρύθμιση της τάσης με το ΣΑΤΥΦ εξετάζεται στο Κεφάλαιο 5. Κάθε
\emph{γραμμή} ΜΤ (feeder) έχει στην αναχώρησή της δικό της διακόπτη ισχύος και ηλεκτρονόμους
υπερέντασης, ενώ στους ζυγούς ΜΤ συνδέονται συστοιχίες πυκνωτών για την αντιστάθμιση της άεργης
ισχύος που απορροφά το δίκτυο διανομής από το σύστημα μεταφοράς.

Οι υποσταθμοί ΜΤ/ΧΤ περιλαμβάνουν έναν μετασχηματιστή 20/0,4 kV ονομαστικής ισχύος από 50 έως
1\,000 kVA, τυπικά με τάση βραχυκύκλωσης 4\% έως τα 630 kVA και 6\% για μεγαλύτερες ισχύες, με
λήψεις που ρυθμίζονται χειροκίνητα μόνο χωρίς φορτίο, συνήθως $\pm 2 \times$ 2,5\%, και την
προστασία του με ασφάλειες ΜΤ ή διακόπτη.

Ο Πίνακας~\ref{tab:hedno} δίνει τα βασικά μεγέθη του ΕΔΔΗΕ. Το μήκος του δικτύου υπερβαίνει τα
250\,000 km, δηλαδή είναι περίπου εικοσαπλάσιο του μήκους των γραμμών του συστήματος μεταφοράς,
και οι υποσταθμοί ΜΤ/ΧΤ πλησιάζουν τους 170\,000, με μέση ονομαστική ισχύ μετασχηματιστή
περίπου 180 kVA. Τα μεγέθη αυτά εξηγούν γιατί η εποπτεία και ο έλεγχος του δικτύου διανομής δεν
μπορούν να στηριχθούν στις μεθόδους του συστήματος μεταφοράς, όπου κάθε ζυγός και κάθε κλάδος
μετρώνται και τηλεχειρίζονται.

\begin{table}[!htbp]
\centering
\begin{tabular}{lr}
\hline
Μέγεθος & Τιμή \\
\hline
Μήκος δικτύου ΜΤ & 121\,267 km \\
Μήκος δικτύου ΧΤ & 132\,194 km \\
Υποσταθμοί ΥΤ/ΜΤ & 246 \\
Υποσταθμοί ΜΤ/ΧΤ & 167\,316 \\
Εγκατεστημένη ισχύς μετασχηματιστών ΜΤ/ΧΤ & 30\,835 MVA \\
Πελάτες ΜΤ & 19\,116 \\
Πελάτες ΧΤ & 7\,655\,055 \\
Διανεμόμενη ενέργεια & 43\,808 GWh \\
Ισχύς σταθμών ΑΠΕ συνδεδεμένων στο δίκτυο & 9\,162 MVA \\
\hline
\end{tabular}
\caption{Βασικά μεγέθη του ΕΔΔΗΕ (2025). Πηγή: ΔΕΔΔΗΕ.}
\label{tab:hedno}
\end{table}

\subsection{Τοπολογία και τρόπος λειτουργίας}
\label{sec:dn-topology}

Τα δίκτυα διανομής σχεδιάζονται με τρεις βασικές τοπολογίες, οι οποίες παρουσιάζονται στο
Σχήμα~\ref{fig:dn-topologies}. Το \emph{ακτινικό} δίκτυο έχει τη μορφή δένδρου: κάθε κόμβος τροφοδοτείται από μία
μόνο διαδρομή από τον υποσταθμό ΥΤ/ΜΤ. Θεωρείται η απλούστερη και οικονομικότερη τοπολογία·
ωστόσο, σε περίπτωση σφάλματος, διακόπτεται η τροφοδότηση όλων των καταναλωτών κατάντη του
στοιχείου προστασίας που ενεργοποιείται. Στο δίκτυο \emph{ανοικτού βρόχου} δύο γραμμές, από
τον ίδιο ή από διαφορετικούς υποσταθμούς, συνδέονται στο άκρο τους μέσω ενός ΚΑ. Το
δίκτυο λειτουργεί ακτινικά, αλλά μετά από ένα σφάλμα το τμήμα που έπαθε βλάβη
απομονώνεται και οι υπόλοιποι χρήστες τροφοδοτούνται από τη γειτονική γραμμή με το
κλείσιμο του ΚΑ. Το \emph{βροχοειδές} δίκτυο λειτουργεί με κλειστούς
βρόχους, οπότε κάθε κόμβος τροφοδοτείται από δύο ή περισσότερες διαδρομές, όπως στο σύστημα
μεταφοράς, και η απώλεια ενός στοιχείου δεν διακόπτει την τροφοδότηση. Η τοπολογία αυτή
απαιτεί όμως πιο σύνθετη προστασία, λειτουργία και έλεγχο, με αποτέλεσμα υψηλότερο κόστος, γι'
αυτό υιοθετείται συνήθως για την τροφοδότηση κρίσιμων φορτίων ή σε κέντρα πόλεων με μεγάλη
πυκνότητα φορτίου.

\begin{figure}[!htbp]
\centering
\includegraphics[width=\linewidth]{figures/dn-topologies.svg}
\caption{Τοπολογίες δικτύων διανομής: ακτινικό δίκτυο, ανοικτός βρόχος και βροχοειδές δίκτυο.}
\label{fig:dn-topologies}
\end{figure}

Η \emph{ακτινική λειτουργία} είναι ο κανόνας στα περισσότερα δίκτυα διανομής. Τα δίκτυα ΜΤ των
αστικών περιοχών κατασκευάζονται ως ανοικτοί βρόχοι και τα εναέρια αγροτικά δίκτυα ως ακτινικά,
με διασυνδέσεις μεταξύ γειτονικών γραμμών μέσω ΚΑ, ενώ τα δίκτυα ΧΤ
είναι σχεδόν πάντοτε ακτινικά. Ο ΔΕΔΔΗΕ περιγράφει το δίκτυο ΜΤ ως βροχοειδές ως προς τη
\emph{δομή} του, όμως οι βρόχοι λειτουργούν ανοικτοί. Η ακτινική λειτουργία απλοποιεί την
προστασία, αφού το ρεύμα σφάλματος ρέει από τον υποσταθμό προς το σφάλμα.
Επιτρέπει επίσης την \emph{αναδιάταξη} (reconfiguration) του δικτύου με αλλαγή της θέσης των ΚΑ,
είτε για την αποκατάσταση της τροφοδότησης μετά από σφάλμα είτε
για τη μείωση των απωλειών και την εξισορρόπηση των φορτίων μεταξύ γραμμών. Η βέλτιστη διάταξη
υπολογίζεται από το σύστημα διαχείρισης της διανομής (Ενότητα «\nameref{sec:adms}»). Η θέση των ΚΑ
δεν είναι επομένως σταθερή, αλλά μεταβάλλεται με τη φόρτιση, τις
εργασίες συντήρησης και τις βλάβες, και ο διαχειριστής πρέπει να γνωρίζει σε κάθε στιγμή την
τρέχουσα τοπολογία. Αξίζει να σημειωθεί ότι η διανεμημένη παραγωγή αναιρεί την παραδοχή της μονόδρομης ροής
ισχύος από τον υποσταθμό προς τα άκρα, με συνέπειες στη ρύθμιση της τάσης και στην προστασία
(Ενότητα «\nameref{sec:active-dn}»).

Η προστασία των ακτινικών γραμμών ΜΤ βασίζεται σε ηλεκτρονόμους υπερέντασης στην αναχώρηση
κάθε γραμμής, σε \emph{διακόπτες αυτόματης επαναφοράς} (reclosers) κατά μήκος των εναέριων
γραμμών, οι οποίοι εκκαθαρίζουν τα παροδικά σφάλματα, την πλειονότητα των σφαλμάτων των
εναέριων γραμμών, με διαδοχικές διακοπές και επαναφορές, και σε ασφάλειες στις διακλαδώσεις
και στους υποσταθμούς ΜΤ/ΧΤ. Οι \emph{διακόπτες φορτίου} χρησιμοποιούνται για την απομόνωση του τμήματος δικτύου με σφάλμα και για την
αναδιάταξη, είτε χειροκίνητα είτε με τηλεχειρισμό (Ενότητα «\nameref{sec:flisr}»).

\subsection{Ιδιαιτερότητες σε σχέση με το σύστημα μεταφοράς}
\label{sec:dn-vs-tn}

Οι μέθοδοι ανάλυσης και λειτουργίας του συστήματος μεταφοράς δεν μεταφέρονται αυτούσιες στα
δίκτυα διανομής, επειδή τα δύο δίκτυα διαφέρουν σε αρκετά σημεία
(Πίνακας~\ref{tab:dn-vs-tn}). Πρώτον, ο λόγος αντίστασης προς αντίδραση $r/x$ των γραμμών
διανομής είναι της τάξης της μονάδας ή μεγαλύτερος, έναντι 0,1 περίπου στις γραμμές των
400 kV. Η ενεργή ισχύς επηρεάζει επομένως την τάση εξίσου ή και περισσότερο από την άεργη,
οπότε η διάκριση μεταξύ του ελέγχου ενεργής ισχύος-συχνότητας και άεργης ισχύος-τάσης, στην
οποία στηρίζεται η λειτουργία του συστήματος μεταφοράς (\hyperref[kef-02-theoria]{Κεφάλαιο 2},
Ενότητα «\nameref{sec:v-q}»), δεν ισχύει, και η ρύθμιση της τάσης με άεργη ισχύ είναι λιγότερο
αποτελεσματική. 

Δεύτερον, η λειτουργία είναι ακτινική, οπότε δεν υπάρχουν παράλληλες διαδρομές ούτε βρόχοι
ροής ισχύος και, χωρίς διανεμημένη παραγωγή, οι τάσεις μειώνονται μονότονα κατά μήκος κάθε
γραμμής. Τρίτον, τα δίκτυα ΧΤ και εν μέρει τα δίκτυα ΜΤ είναι \emph{ασύμμετρα}: τα
περισσότερα φορτία και πολλές μονάδες παραγωγής είναι μονοφασικά και κατανέμονται άνισα στις
τρεις φάσεις, οπότε το μονοφασικό ισοδύναμο κύκλωμα, που αρκεί για το σύστημα μεταφοράς, δεν
είναι πάντοτε επαρκές (Ενότητα «\nameref{sec:dsse}» και Κεφάλαιο 5).
Τέταρτον, το πλήθος των κόμβων είναι τεράστιο, ενώ οι μετρήσεις σε πραγματικό χρόνο είναι
ελάχιστες και περιορίζονται κατά κανόνα στον υποσταθμό ΥΤ/ΜΤ· τα φορτία των υποσταθμών
ΜΤ/ΧΤ εκτιμώνται από στατιστικά προφίλ και όχι από μετρήσεις. Τέλος, οι διαταραχές που
απασχολούν τον διαχειριστή διανομής είναι κυρίως οι διακοπές της τροφοδότησης, οι αποκλίσεις
του μέτρου της τάσης και η υπερφόρτιση των στοιχείων, και όχι η ευστάθεια γωνίας ή
συχνότητας, η οποία αφορά το ΣΜΗΕ (\hyperref[kef-03-theoria]{Κεφάλαιο 3}).

\begin{table}[!htbp]
\centering
\begin{tabular}{lll}
\hline
Χαρακτηριστικό & Σύστημα μεταφοράς & Δίκτυο διανομής \\
\hline
Επίπεδα τάσης & 400, 150 kV & 20 kV, 400 V \\
Τοπολογία & βροχοειδής & ακτινική (ανοικτοί βρόχοι) \\
Ροή ισχύος & αμφίδρομη, ελεγχόμενη & υποσταθμός $\rightarrow$ φορτία, αντίστροφη με διανεμημένη παραγωγή \\
Λόγος $r/x$ γραμμών & 0,05--0,3 & 0,5--3 \\
Συμμετρία φάσεων & σχεδόν συμμετρικό & ασύμμετρα φορτία και παραγωγή \\
Ζυγοί ή κόμβοι & εκατοντάδες & εκατοντάδες χιλιάδες \\
Μετρήσεις πραγματικού χρόνου & σε κάθε ζυγό και κλάδο & κυρίως στον υποσταθμό ΥΤ/ΜΤ \\
Μοντέλο φορτίου & μετρούμενες εγχύσεις & στατιστικά προφίλ, έξυπνοι μετρητές \\
Κύριες διαταραχές & ευστάθεια, συμφόρηση & διακοπές, αποκλίσεις τάσης, υπερφόρτιση \\
\hline
\end{tabular}
\caption{Σύγκριση των χαρακτηριστικών του συστήματος μεταφοράς και των δικτύων διανομής.}
\label{tab:dn-vs-tn}
\end{table}

\section{Απαιτήσεις λειτουργίας}
\label{sec:dn-requirements}

Οι απαιτήσεις λειτουργίας του δικτύου διανομής απορρέουν από το ευρωπαϊκό και το εθνικό
ρυθμιστικό πλαίσιο. Η Οδηγία (ΕΕ) 2019/944, η οποία ενσωματώθηκε με τον ν. 5037/2023,
αναθέτει στον \emph{διαχειριστή δικτύου διανομής} (Distribution System Operator, DSO) τη
λειτουργία, τη συντήρηση και την ανάπτυξη ασφαλούς, αξιόπιστου και αποδοτικού δικτύου, τη
σύνδεση των χρηστών χωρίς διακρίσεις, τη διαχείριση των δεδομένων μέτρησης και την προμήθεια
υπηρεσιών ευελιξίας από ΔΕΠ, ενώ ο Κώδικας Διαχείρισης του ΕΔΔΗΕ, τον οποίο
εγκρίνει η ΡΑΑΕΥ, εξειδικεύει τους τεχνικούς κανόνες. Οι δύο βασικές απαιτήσεις προς τους
χρήστες είναι η ποιότητα της τάσης και η συνέχεια της τροφοδότησης.

\subsection{Χαρακτηριστικά της τάσης: το πρότυπο EN 50160}
\label{sec:en50160}

Η ποιότητα της τάσης που παραδίδεται στους χρήστες των δικτύων διανομής ορίζεται από το
ευρωπαϊκό πρότυπο EN 50160 «Χαρακτηριστικά της τάσης που παρέχεται από τα δημόσια δίκτυα
διανομής», το οποίο εφαρμόζει ο ΔΕΔΔΗΕ ως ελληνικό πρότυπο ΕΛΟΤ EN 50160. Το
πρότυπο περιγράφει τα χαρακτηριστικά της τάσης στο σημείο παράδοσης υπό κανονικές συνθήκες
λειτουργίας, δηλαδή εκτός σφαλμάτων και έκτακτων καταστάσεων
(Πίνακας~\ref{tab:en50160}). Κάθε μέγεθος, π.χ. η ενεργός τιμή της τάσης, δεν ελέγχεται
στιγμιαία, αλλά ως μέση τιμή σε ένα σταθερό χρονικό διάστημα (διάστημα ανάλυσης), π.χ. 10
λεπτά. Οι μέσες τιμές συγκεντρώνονται σε μια περίοδο παρατήρησης, π.χ. μίας εβδομάδας, και το
πρότυπο απαιτεί να βρίσκεται εντός των ορίων ένα ελάχιστο ποσοστό τους, π.χ. 95\%.

\begin{table}[!htbp]
\centering
\begin{tabular}{llll}
\hline
Χαρακτηριστικό & Μέγεθος & Όριο & Ποσοστό μετρήσεων \\
\hline
\multirow{2}{*}{Συχνότητα} & \multirow{2}{*}{μέση τιμή 10 s} & 50 Hz $\pm 1\%$ & 99,5\% του έτους \\
 & & $+4\%$ / $-6\%$ & 100\% \\
\multirow{2}{*}{Τάση ΧΤ} & \multirow{2}{*}{μέση rms 10 min} & $U_n \pm 10\%$ & 95\% της εβδομάδας \\
 & & $U_n$ $+10\%$ / $-15\%$ & 100\% \\
\multirow{2}{*}{Τάση ΜΤ} & \multirow{2}{*}{μέση rms 10 min} & $U_n \pm 10\%$ & 99\% της εβδομάδας \\
 & & $U_n \pm 15\%$ & 100\% \\
Ασυμμετρία (ΧΤ και ΜΤ) & $u_2$, 10 min & 2\% & 95\% της εβδομάδας \\
Αρμονική παραμόρφωση & THD, 10 min & 8\% & 95\% της εβδομάδας \\
\hline
\end{tabular}
\caption{Κύρια χαρακτηριστικά της τάσης κατά το EN 50160 για διασυνδεδεμένα συστήματα.}
\label{tab:en50160}
\end{table}

Στον Πίνακα~\ref{tab:en50160}, $U_n$ είναι η ονομαστική τάση του δικτύου. Στη ΧΤ, το 95\% των
μέσων τιμών 10 λεπτών μιας εβδομάδας πρέπει να βρίσκεται εντός $\pm 10\%$ της ονομαστικής
τάσης και όλες εντός $+10\%$ / $-15\%$. Στη ΜΤ, το 99\% των τιμών πρέπει να βρίσκεται εντός
$\pm 10\%$ και όλες εντός $\pm 15\%$ της ονομαστικής τάσης. Ο διαχειριστής επιμερίζει το
διαθέσιμο εύρος $\pm 10\%$ μεταξύ των δικτύων ΜΤ και ΧΤ και επιλέγει τα μέσα ρύθμισης ώστε οι
τάσεις όλων των χρηστών να παραμένουν εντός των ορίων ανεξάρτητα από τις συνθήκες φόρτισης του
δικτύου (Κεφάλαιο 5).

Η \emph{ασυμμετρία} της τάσης, η οποία προκαλείται κυρίως από την άνιση κατανομή των μονοφασικών
φορτίων και μονάδων ΔΕΠ στις φάσεις, ορίζεται από τον λόγο $u_2 = |\bar{V}_2|/|\bar{V}_1|$
της αντίστροφης προς την ορθή συνιστώσα, ο οποίος δεν πρέπει να υπερβαίνει το 2\% για το 95\%
των μετρήσεων της εβδομάδας. Για την ομοπολική συνιστώσα της τάσης το πρότυπο δεν θέτει
περιορισμό. Ο λόγος υπολογίζεται από τους φασιθέτες των θεμελιωδών
συνιστωσών με τις συμμετρικές συνιστώσες (Κεφάλαιο 5) ή, απλούστερα, μόνο
από τα μέτρα των τριών πολικών τάσεων $U_{12}$, $U_{23}$, $U_{31}$:

\begin{equation}
u_2 = \sqrt{\frac{1 - \sqrt{3 - 6\beta}}{1 + \sqrt{3 - 6\beta}}}, \qquad
\beta = \frac{U_{12}^4 + U_{23}^4 + U_{31}^4}{\left(U_{12}^2 + U_{23}^2 + U_{31}^2\right)^2}
\label{eq:vuf-line}
\end{equation}

Για συμμετρικές τάσεις $\beta = 1/3$ και $u_2 = 0$. Η \emph{αρμονική παραμόρφωση} περιορίζεται
τόσο ανά αρμονική, με χωριστό όριο για κάθε τάξη $h$, όσο και συνολικά, με τον συντελεστή
ολικής αρμονικής παραμόρφωσης

\begin{equation}
\mathrm{THD} = \sqrt{\sum_{h=2}^{40} u_h^2}
\label{eq:thd}
\end{equation}

όπου $u_h$ το πλάτος της αρμονικής τάξης $h$ ως ποσοστό της θεμελιώδους. Ο THD δεν πρέπει να
υπερβαίνει το 8\% για το 95\% των μετρήσεων της εβδομάδας, τόσο στη ΧΤ όσο και στη ΜΤ. Το
πρότυπο καλύπτει επίσης τις βυθίσεις και τις υπερτάσεις της τάσης, τις διακοπές και τις
τάσεις σηματοδότησης· για τις βυθίσεις, τις υπερτάσεις και τις διακοπές δίνει μόνο τρόπους
ταξινόμησης και ενδεικτικές τιμές, όχι όρια.

\subsection{Αξιοπιστία τροφοδότησης}
\label{sec:reliability}

Η δεύτερη βασική απαίτηση είναι η \emph{συνέχεια της τροφοδότησης}. Επειδή τα δίκτυα
διανομής λειτουργούν ακτινικά, κάθε σφάλμα προκαλεί διακοπή σε ένα σύνολο καταναλωτών, και
η πλειονότητα των διακοπών που αντιλαμβάνονται οι καταναλωτές οφείλεται στο δίκτυο διανομής
και όχι στο σύστημα μεταφοράς ή στην παραγωγή. Η αξιοπιστία μετράται με δείκτες που
συνοψίζουν τις \emph{μακράς διάρκειας} διακοπές ενός έτους, δηλαδή τις διακοπές άνω των
τριών λεπτών κατά το EN 50160. Οι πιο διαδεδομένοι είναι ο δείκτης μέσης διάρκειας διακοπών
του συστήματος (System Average Interruption Duration Index, SAIDI), ο δείκτης μέσης συχνότητας
διακοπών του συστήματος (System Average Interruption Frequency Index, SAIFI) και ο δείκτης
μέσης διάρκειας διακοπής καταναλωτή (Customer Average Interruption Duration Index, CAIDI):

\begin{equation}
\mathrm{SAIDI} = \frac{\sum_i N_i\,t_i}{N_T}, \qquad
\mathrm{SAIFI} = \frac{\sum_i N_i}{N_T}, \qquad
\mathrm{CAIDI} = \frac{\mathrm{SAIDI}}{\mathrm{SAIFI}} = \frac{\sum_i N_i\,t_i}{\sum_i N_i}
\label{eq:saidi}
\end{equation}

όπου το άθροισμα εκτείνεται σε όλες τις διακοπές $i$ του έτους, $N_i$ το πλήθος των
καταναλωτών που διακόπηκαν στη διακοπή $i$, $t_i$ η διάρκειά της και $N_T$ το συνολικό πλήθος
των καταναλωτών. Ο SAIDI εκφράζεται σε λεπτά ανά καταναλωτή και έτος, ο SAIFI σε διακοπές
ανά καταναλωτή και έτος και ο CAIDI, η μέση διάρκεια μιας διακοπής, σε λεπτά. Αξίζει να
σημειωθεί ότι οι διακοπές βραχείας διάρκειας, που οφείλονται κυρίως στην αυτόματη επαναφορά των
διακοπτών, μετρώνται χωριστά με τον δείκτη MAIFI (Momentary Average Interruption
Frequency Index), ο οποίος ορίζεται όπως ο SAIFI. Όταν το πλήθος των καταναλωτών κάθε διακοπής
δεν είναι γνωστό, οι δείκτες σταθμίζονται με την ισχύ που διακόπηκε:

\begin{equation}
\mathrm{ASIDI} = \frac{\sum_i L_i\,t_i}{L_T}, \qquad
\mathrm{ASIFI} = \frac{\sum_i L_i}{L_T}, \qquad
\mathrm{ENS} = \sum_i P_i\,t_i
\label{eq:asidi}
\end{equation}

όπου $L_i$ η ονομαστική ισχύς των μετασχηματιστών ΜΤ/ΧΤ που διακόπηκαν στη διακοπή $i$, $L_T$
η συνολική ισχύς των μετασχηματιστών και $P_i$ το φορτίο που δεν εξυπηρετήθηκε· η μη
διανεμόμενη ενέργεια (Energy Not Supplied, ENS), σε MWh, χρησιμοποιείται για την αποτίμηση
του κόστους των διακοπών και στα ρυθμιστικά κίνητρα. Οι δείκτες υπολογίζονται χωριστά για τις
προγραμματισμένες διακοπές, για τις οποίες οι χρήστες ειδοποιούνται εκ των προτέρων, και για
τις μη προγραμματισμένες. Επειδή λίγες ημέρες με καταιγίδες, πυρκαγιές ή χιονοπτώσεις
μπορούν να καθορίσουν τον ετήσιο SAIDI, οι μη προγραμματισμένες διακοπές αναφέρονται και
χωρίς τα \emph{εξαιρετικά συμβάντα}, τα οποία κάθε χώρα ορίζει με τους δικούς της κανόνες.

\begin{figure}[!htbp]
\centering
\includegraphics[width=\linewidth]{figures/ceer-saidi.svg}
\caption{SAIDI και SAIFI των μη προγραμματισμένων διακοπών μακράς διάρκειας ανά χώρα το 2018.
Πηγή: CEER, 7η συγκριτική έκθεση (2022).}
\label{fig:ceer-saidi}
\end{figure}

Το Συμβούλιο Ευρωπαϊκών Ρυθμιστικών Αρχών Ενέργειας (CEER) δημοσιεύει περιοδικά συγκριτικές
εκθέσεις για την ποιότητα της τροφοδότησης στην Ευρώπη. Το Σχήμα~\ref{fig:ceer-saidi} δείχνει
τον SAIDI (πάνω) και τον SAIFI (κάτω) των μη προγραμματισμένων διακοπών μακράς διάρκειας κατά
την τελευταία έκθεση. Στην Ελλάδα, το 2018, ο SAIDI ήταν 119 min χωρίς τα εξαιρετικά
συμβάντα και 173 min με όλα τα συμβάντα, και ο SAIFI 1,62 και 2,33 διακοπές αντίστοιχα,
έναντι 14 min και 0,25 διακοπών στη Γερμανία.

Η διάρκεια μιας διακοπής αναλύεται σε τρία στάδια: τον χρόνο εντοπισμού του σφάλματος, τον
χρόνο απομόνωσης του τμήματος που έπαθε βλάβη και επανατροφοδότησης των υπόλοιπων
καταναλωτών, και τον χρόνο επισκευής. Στα συμβατικά δίκτυα τα δύο πρώτα στάδια απαιτούν τη
μετακίνηση συνεργείου και χειρισμούς επί τόπου και διαρκούν από μισή έως δύο ώρες· η
υποδομή τηλεχειρισμού και οι αυτοματισμοί της Ενότητας «\nameref{sec:flisr}» μπορούν να τα
περιορίσουν σε λίγα λεπτά, μειώνοντας έτσι τον SAIDI χωρίς να αλλάζει το πλήθος των
σφαλμάτων. Η μείωση του SAIFI απαιτεί αντίθετα επεμβάσεις στο ίδιο το δίκτυο, όπως
υπογειοποίηση, μονωμένους αγωγούς και κλαδέματα, ή μικρότερες ζώνες προστασίας.

\section{Από το παθητικό στο ενεργό δίκτυο}
\label{sec:active-dn}

Οι \emph{διανεμημένοι ενεργειακοί πόροι} (Distributed Energy Resources, DER· ΔΕΠ) είναι οι
μονάδες παραγωγής, αποθήκευσης και ελεγχόμενης κατανάλωσης που συνδέονται στα δίκτυα
διανομής, όπως φωτοβολταϊκοί σταθμοί, μικρά αιολικά πάρκα, μικροί υδροηλεκτρικοί σταθμοί,
συσσωρευτές, φορτιστές ηλεκτρικών οχημάτων, αντλίες θερμότητας και ευέλικτα φορτία. Η
\emph{διανεμημένη παραγωγή} είναι το μέρος των ΔΕΠ που παράγει ενέργεια. Στο ΕΔΔΗΕ η ισχύς των
σταθμών ΑΠΕ, κυρίως φωτοβολταϊκών, ξεπερνούσε το 2025 τα 9 GVA, περίπου το μισό της
εγκατεστημένης ισχύος ΑΠΕ της χώρας χωρίς τους μεγάλους υδροηλεκτρικούς σταθμούς.

Η διείσδυση αυτή μεταβάλλει τη λειτουργία του δικτύου με αρκετούς τρόπους. Η ροή ισχύος
αντιστρέφεται στις γραμμές και στους μετασχηματιστές, οπότε η τάση ανυψώνεται κατά μήκος των
γραμμών αντί να μειώνεται και γίνεται μέγιστη τις ώρες της μεγάλης παραγωγής και
της μικρής ζήτησης. Η ανύψωση στο σημείο σύνδεσης μιας μονάδας είναι, όπως στην
\eqref{eq:dv-ssc} του Κεφαλαίου 2, αντιστρόφως ανάλογη της ισχύος βραχυκύκλωσης του σημείου,
αλλά, λόγω της σημαντικής αντίστασης των γραμμών, την προκαλεί κυρίως η ενεργός ισχύς της
μονάδας, ενώ η απορρόφηση άεργης ισχύος τη μειώνει. Οι γραμμές και οι μετασχηματιστές
φορτίζονται από την παραγωγή και όχι μόνο από τη ζήτηση, για την οποία διαστασιολογήθηκαν. Το
ρεύμα βραχυκύκλωσης αυξάνεται, επειδή κάθε μονάδα συνεισφέρει στο σφάλμα, ενώ η συμβολή των
μετατροπέων διαφέρει από εκείνη των σύγχρονων μηχανών, οπότε η ρύθμιση των προστασιών
αναθεωρείται.

Για την αντιμετώπιση των παραπάνω προβλημάτων εισάγεται η έννοια του \emph{έξυπνου δικτύου}
(smart grid): ένα δίκτυο στο οποίο η παραγωγή, η αποθήκευση και τα ευέλικτα φορτία
συμμετέχουν στη λειτουργία, με τοπικούς ελέγχους ή με εντολές του διαχειριστή, και στο οποίο ο
διαχειριστής διαθέτει μετρήσεις και μέσα ελέγχου, ώστε να εξασφαλίσει αυξημένα ποσοστά
διείσδυσης ΔΕΠ. Η ενεργή αυτή διαχείριση αντιπαρατίθεται στην παραδοσιακή πρακτική κατά την
οποία ο σταθμός, μετά τη σύνδεσή του, λειτουργεί χωρίς καμία παρακολούθηση ή ρύθμιση από τον
διαχειριστή (fit and forget). Οι τεχνικές ρύθμισης της τάσης και της συμφόρησης που υλοποιούν
την ενεργή διαχείριση εξετάζονται στα επόμενα κεφάλαια· προϋπόθεσή τους είναι η υποδομή
εποπτείας και ελέγχου που ακολουθεί.

\section{Εποπτεία και έλεγχος}
\label{sec:dn-control}

Η υποδομή εποπτείας και ελέγχου ενός δικτύου διανομής έχει τρία επίπεδα (Σχήμα~\ref{fig:adms}):
τις συσκευές πεδίου, δηλαδή τους μετρητές, τους αισθητήρες, τις ευφυείς ηλεκτρονικές
συσκευές και τους τηλεχειριζόμενους διακόπτες, το τηλεπικοινωνιακό δίκτυο που τις συνδέει με
τα κέντρα ελέγχου, και τα συστήματα των κέντρων ελέγχου, δηλαδή το SCADA και τις εφαρμογές
διαχείρισης της διανομής. Η αρχιτεκτονική είναι η ίδια με του συστήματος μεταφοράς
(Κεφάλαιο 1, Ενότητα «\nameref{sec:meas-chain}»), αλλά η κλίμακα διαφέρει κατά δύο έως τρεις
τάξεις μεγέθους.

\subsection{Έξυπνοι μετρητές και προηγμένη υποδομή μέτρησης}
\label{sec:smart-meters}

Ο \emph{έξυπνος μετρητής} (smart meter) είναι ηλεκτρονικός μετρητής ενέργειας με
αμφίδρομη επικοινωνία (Σχήμα~\ref{fig:smart-meter-photo}). Εκτός από την ενέργεια που καταναλώνεται και, για τους
αυτοπαραγωγούς, παράγεται, καταγράφει την ενεργή και την άεργη ενέργεια ανά διάστημα
ολοκλήρωσης, συνήθως 15 λεπτών, τη μέση τάση κάθε φάσης και συμβάντα ποιότητας όπως
διακοπές, βυθίσεις και υπερτάσεις, και επιτρέπει την απομακρυσμένη ανάγνωση, τη
σύνδεση και αποσύνδεση της παροχής και τον περιορισμό της ισχύος. Οι μετρήσεις
μεταδίδονται στον διαχειριστή μέσω της \emph{προηγμένης υποδομής μέτρησης} (Advanced
Metering Infrastructure, AMI), η οποία περιλαμβάνει τους μετρητές, το δίκτυο επικοινωνίας,
το κεντρικό σύστημα συλλογής (Head-End System, HES) που επικοινωνεί με τους μετρητές, και
το σύστημα διαχείρισης δεδομένων μέτρησης (Meter Data Management System, MDMS), το οποίο
επικυρώνει, αποθηκεύει και διαθέτει τα δεδομένα στην τιμολόγηση, στους προμηθευτές, στους
χρήστες και στις εφαρμογές λειτουργίας του δικτύου (Σχήμα~\ref{fig:ami}). Η ανταλλαγή των
δεδομένων ακολουθεί το πρότυπο DLMS/COSEM της IEC 62056.

\begin{figure}[!htbp]
\centering
\includegraphics[width=0.36\linewidth]{figures/smart-meter-photo.jpg}
\caption{Έξυπνος μετρητής με μονάδα επικοινωνίας. Φωτογραφία: Fredrik Lindseth,
\href{https://commons.wikimedia.org/wiki/File:Kaifa_MA105H2E_AMS-måler.jpg}{«Kaifa MA105H2E
AMS-måler»}, άδεια \href{https://creativecommons.org/licenses/by-sa/4.0/}{CC BY-SA 4.0}, μέσω
Wikimedia Commons.}
\label{fig:smart-meter-photo}
\end{figure}

\begin{figure}[!htbp]
\centering
\includegraphics[width=\linewidth]{figures/ami.svg}
\caption{Προηγμένη υποδομή μέτρησης.}
\label{fig:ami}
\end{figure}

Η επικοινωνία των μετρητών με το HES γίνεται με δύο αρχιτεκτονικές. Στην
\emph{άμεση} αρχιτεκτονική κάθε μετρητής διαθέτει μόντεμ κινητής τηλεφωνίας, 4G ή δικτύων
χαμηλής ισχύος όπως τα NB-IoT και LTE-M, και επικοινωνεί απευθείας με το HES. Στην
\emph{ιεραρχική} αρχιτεκτονική οι μετρητές μιας γραμμής ΧΤ επικοινωνούν με έναν
\emph{συγκεντρωτή δεδομένων} (data concentrator) στον υποσταθμό ΜΤ/ΧΤ, είτε με
ραδιοσυχνότητες είτε μέσω των ίδιων των αγωγών ΧΤ (Power Line Communication, PLC), και ο
συγκεντρωτής με το HES.

Ο όγκος των δεδομένων ενός συστήματος AMI είναι

\begin{equation}
D = N_m\,n_c\,\frac{T}{\Delta t}\,b
\label{eq:ami-data}
\end{equation}

όπου $N_m$ το πλήθος των μετρητών, $n_c$ τα κανάλια μέτρησης ανά μετρητή, $\Delta t$ το
διάστημα ολοκλήρωσης, $T$ η περίοδος και $b$ τα byte ανά τιμή. Για 7,7 εκατομμύρια μετρητές
με πέντε κανάλια, όπως σε έναν μονοφασικό μετρητή η ενεργή και η άεργη ενέργεια στις δύο
κατευθύνσεις και η μέση τάση, ανά 15 λεπτά ($T/\Delta t = 96$ ανά ημέρα) και 8 byte ανά τιμή
με τη χρονοσφραγίδα και την ένδειξη ποιότητας, προκύπτουν περίπου 30 GB ημερησίως ή 11 TB
ετησίως μόνο για τις καμπύλες φορτίου, χωρίς τα συμβάντα, τα μεταδεδομένα και τα αντίγραφα
ασφαλείας.

Στην Ελλάδα, οι μετρητές καταγράφουν την ενέργεια ανά 15 λεπτά και επικοινωνούν με το HES μέσω
κινητής τηλεφωνίας, ενώ το MDMS σχεδιάζεται για περίπου 7,7 εκατομμύρια σημεία μέτρησης. Για τον διαχειριστή, πέρα από την τιμολόγηση και την
αντιμετώπιση των μη τεχνικών απωλειών, το κύριο όφελος είναι η \emph{παρατηρησιμότητα των
δικτύων ΧΤ}, τα οποία χωρίς τους έξυπνους μετρητές
δεν διαθέτουν καμία μέτρηση πέρα από τον υποσταθμό ΜΤ/ΧΤ, και η άμεση γνώση των διακοπών
τροφοδότησης από τα μηνύματα που στέλνουν οι μετρητές όταν χάνουν την τάση, με την ενέργεια
που έχουν αποθηκεύσει σε πυκνωτές (last gasp). Επειδή μέρος των μηνυμάτων αυτών μπορεί να
χαθεί, π.χ. όταν διακόπτονται ταυτόχρονα πολλοί μετρητές, οι μετρητές καταγράφουν επιπλέον την
έναρξη και τη λήξη κάθε διακοπής στη μνήμη τους. Οι καταγραφές αυτές βελτιώνουν και την
ακρίβεια υπολογισμού των δεικτών αξιοπιστίας της \eqref{eq:saidi}, όπως ο SAIDI και ο SAIFI,
επειδή το πλήθος $N_i$ των καταναλωτών και η διάρκεια $t_i$ κάθε διακοπής προκύπτουν πλέον από
μετρήσεις και όχι από εκτιμήσεις ή από τη στάθμιση με την ισχύ της \eqref{eq:asidi}.

\subsection{Ευφυείς ηλεκτρονικές συσκευές και ψηφιακοί υποσταθμοί}
\label{sec:ied}

Οι \emph{ευφυείς ηλεκτρονικές συσκευές} (Intelligent Electronic Devices, IED) είναι συσκευές
με μικροεπεξεργαστή που δειγματοληπτούν τις τάσεις και τα ρεύματα των μετασχηματιστών
μέτρησης, υπολογίζουν τα παράγωγα μεγέθη, λαμβάνουν αποφάσεις τοπικά και επικοινωνούν με το
κέντρο ελέγχου και μεταξύ τους. Στα δίκτυα διανομής οι συνηθέστερες είναι οι ηλεκτρονόμοι
προστασίας των γραμμών και των μετασχηματιστών, οι ελεγκτές των ΣΑΤΥΦ, οι οποίοι υλοποιούν
την αυτόματη ρύθμιση της τάσης των ζυγών ΜΤ (Κεφάλαιο 5), οι
ελεγκτές των διακοπτών αυτόματης επαναφοράς και των τηλεχειριζόμενων διακοπτών φορτίου, οι
ελεγκτές των συστοιχιών πυκνωτών, κτλ. Μια σύγχρονη IED προστασίας
(Σχήμα~\ref{fig:ied-photo}) συνδυάζει σε μία συσκευή πολλές λειτουργίες προστασίας, μέτρηση,
καταγραφή συμβάντων και κυματομορφών διαταραχών, ελέγχους διακοπτών και επικοινωνία με
πρωτόκολλα όπως το IEC 61850 εντός του υποσταθμού και τα IEC 60870-5-104 ή DNP3 προς το
κέντρο ελέγχου (Κεφάλαιο 1, Ενότητα «\nameref{sec:meas-chain}»).

\begin{figure}[!htbp]
\centering
\includegraphics[width=0.36\linewidth]{figures/ied-photo.jpg}
\caption{IED προστασίας και ελέγχου πεδίου ΜΤ, με το μονογραμμικό διάγραμμα του πεδίου στην
οθόνη. Φωτογραφία: Zen 38, \href{https://commons.wikimedia.org/wiki/File:Protective_relay.jpg}{«Protective
relay»}, άδεια \href{https://creativecommons.org/licenses/by-sa/3.0/}{CC BY-SA 3.0}, μέσω
Wikimedia Commons.}
\label{fig:ied-photo}
\end{figure}

Στον \emph{συμβατικό υποσταθμό} κάθε σήμα μεταφέρεται με χωριστό καλώδιο χαλκού
(Σχήμα~\ref{fig:substation-cabling}): τα δευτερεύοντα των μετασχηματιστών έντασης και τάσης
καλωδιώνονται από το πεδίο, μέσω των καναλιών καλωδίων, προς κάθε ηλεκτρονόμο και όργανο
μέτρησης στο κτίριο ελέγχου, και κάθε εντολή και ένδειξη θέσης διακόπτη ή αποζεύκτη
μεταφέρεται με δικό της ζεύγος αγωγών. Σε έναν υποσταθμό ΥΤ/ΜΤ τα καλώδια αυτά φθάνουν
συνολικά σε μήκος πολλών χιλιομέτρων. Η καλωδίωση είναι ακριβή, κάθε αλλαγή απαιτεί επέμβαση
στο πεδίο και νέες δοκιμές. Στον \emph{ψηφιακό υποσταθμό} (digital substation) τα αναλογικά
σήματα ψηφιοποιούνται κοντά στην πηγή τους και μεταφέρονται μέσω λίγων οπτικών ινών σε κοινό
δίκτυο επικοινωνίας, σύμφωνα με το πρότυπο IEC 61850.

\begin{figure}[!htbp]
\centering
\includegraphics[width=\linewidth]{figures/substation-cabling.svg}
\caption{Καλωδίωση συμβατικού (πάνω) και ψηφιακού (κάτω) υποσταθμού.}
\label{fig:substation-cabling}
\end{figure}

Το IEC 61850 ορίζει ένα κοινό μοντέλο δεδομένων, στο οποίο κάθε λειτουργία, π.χ. ένας
διακόπτης ισχύος, μια προστασία υπερέντασης ή μια μέτρηση, παριστάνεται ως \emph{λογικός
κόμβος} (logical node) με τυποποιημένο όνομα και δεδομένα (Πίνακας~\ref{tab:iec61850-ln}),
ώστε συσκευές διαφορετικών κατασκευαστών να συνεργάζονται χωρίς ειδικές προσαρμογές. Ορίζει
επίσης τρεις υπηρεσίες επικοινωνίας:

\begin{table}[!htbp]
\centering
\begin{tabular}{lll}
\hline
Ομάδα & Παράδειγμα & Λειτουργία \\
\hline
P (προστασία) & PTOC & προστασία υπερέντασης \\
C (έλεγχος) & CSWI & έλεγχος διακόπτη \\
X (διακοπτικός εξοπλισμός) & XCBR & διακόπτης ισχύος \\
M (μέτρηση) & MMXU & τριφασικές μετρήσεις \\
A (αυτόματος έλεγχος) & ATCC & αυτόματος έλεγχος ΣΑΤΥΦ \\
\hline
\end{tabular}
\caption{Παραδείγματα λογικών κόμβων του IEC 61850.}
\label{tab:iec61850-ln}
\end{table}

\begin{itemize}
\item \textbf{Δειγματοληπτημένες τιμές} (Sampled Values, SV), κατά το IEC 61850-9-2. Οι
      \emph{μονάδες συγχώνευσης} (merging units) δειγματοληπτούν τα δευτερεύοντα των
      μετασχηματιστών μέτρησης και δημοσιεύουν τα δείγματα στον \emph{δίαυλο διεργασίας}
      (process bus), από τον οποίο τα λαμβάνουν όλες οι IED που τα χρειάζονται.
\item \textbf{Μηνύματα GOOSE} (Generic Object Oriented Substation Event), κατά το
      IEC 61850-8-1. Μηνύματα από συσκευή σε συσκευή, με χρόνο μετάδοσης λίγων ms, που
      μεταφέρουν εντολές ζεύξης και μανδαλώσεις και αντικαθιστούν τις καλωδιώσεις μεταξύ των
      ηλεκτρονόμων και προς τους διακόπτες.
\item \textbf{Επικοινωνία πελάτη-εξυπηρετητή} κατά το πρότυπο MMS, για την ανταλλαγή
      μετρήσεων, ενδείξεων και εντολών μέσω του \emph{διαύλου σταθμού} (station bus) και, μέσω
      της πύλης επικοινωνίας του υποσταθμού, με το κέντρο ελέγχου.
\end{itemize}

\begin{figure}[!htbp]
\centering
\includegraphics[width=\linewidth]{figures/digital-substation.svg}
\caption{Συμβατικός (αριστερά) και ψηφιακός (δεξιά) υποσταθμός σε ένα πεδίο αναχώρησης γραμμής
ΜΤ.}
\label{fig:digital-substation}
\end{figure}

Το Σχήμα~\ref{fig:digital-substation} συγκρίνει τις δύο λύσεις σε ένα πεδίο αναχώρησης γραμμής
ΜΤ, στα τρία επίπεδα του υποσταθμού: διεργασίας, πεδίου και σταθμού. Στον συμβατικό υποσταθμό
τα δευτερεύοντα των μετασχηματιστών μέτρησης, οι ενδείξεις και οι εντολές του διακόπτη φθάνουν
με χωριστά καλώδια σε κάθε ηλεκτρονόμο, όργανο και στη \emph{μονάδα τηλεχειρισμού}
(Remote Terminal Unit, RTU), η οποία συγκεντρώνει τα δεδομένα του υποσταθμού, τα μεταδίδει
στο κέντρο ελέγχου και εκτελεί τις εντολές τηλεχειρισμού (Κεφάλαιο 1, Ενότητα
«\nameref{sec:meas-chain}»). Στον ψηφιακό υποσταθμό η μονάδα συγχώνευσης δημοσιεύει τις SV
στον δίαυλο διεργασίας, οι IED στέλνουν τις εντολές ζεύξης (μηνύματα
GOOSE) στην IED του διακόπτη και επικοινωνούν με την πύλη μέσω του διαύλου σταθμού, ενώ όλες
οι συσκευές συγχρονίζονται με ακρίβεια μικροδευτερολέπτου από ένα κύριο ρολόι με το
πρωτόκολλο ακριβούς χρόνου (Precision Time Protocol, PTP), ώστε τα δείγματα διαφορετικών
μονάδων να συνδυάζονται σωστά, π.χ. στη διαφορική προστασία.

\subsection{Τηλεπικοινωνίες}
\label{sec:dn-comms}

Το τηλεπικοινωνιακό δίκτυο ενός διαχειριστή διανομής συνδυάζει αρκετές τεχνολογίες, ανάλογα
με τις απαιτήσεις κάθε εφαρμογής σε καθυστέρηση, αξιοπιστία, όγκο δεδομένων και κόστος
(Πίνακας~\ref{tab:comm-req}). Οι υποσταθμοί ΥΤ/ΜΤ και τα κέντρα ελέγχου συνδέονται με οπτικές
ίνες με μικρή καθυστέρηση και μεγάλη χωρητικότητα. Οι τηλεχειριζόμενοι διακόπτες και οι
διακόπτες αυτόματης επαναφοράς των γραμμών ΜΤ, οι υποσταθμοί ΜΤ/ΧΤ, οι σταθμοί ΑΠΕ και οι
έξυπνοι μετρητές επικοινωνούν κυρίως μέσω των δημόσιων δικτύων κινητής τηλεφωνίας (4G/5G).
Κοινές απαιτήσεις όλων των επικοινωνιών είναι η διαλειτουργικότητα συσκευών διαφορετικών κατασκευαστών, την οποία εξασφαλίζουν ανοιχτά
πρότυπα όπως τα IEC 61850, IEC 60870-5-104 και DLMS/COSEM, και η κυβερνοασφάλεια
(Κεφάλαιο 8).

\begin{table}[!htbp]
\centering
\begin{tabular}{llll}
\hline
Εφαρμογή & Καθυστέρηση & Όγκος δεδομένων & Τυπικό μέσο \\
\hline
Προστασία με GOOSE, SV & $\le$ 3 ms & Mbit/s ανά μονάδα & Ethernet υποσταθμού \\
Διαφορική προστασία γραμμής & 5--10 ms & 64 kbit/s έως 2 Mbit/s & οπτικές ίνες \\
SCADA υποσταθμών ΥΤ/ΜΤ & 1--2 s & kbit/s & οπτικές ίνες \\
Τηλεχειρισμοί, FLISR & δευτερόλεπτα & kbit/s & 4G/5G, ραδιοδίκτυα \\
Εντολές σε σταθμούς ΑΠΕ & δευτερόλεπτα & kbit/s & 4G/5G \\
Συμβάντα μετρητών & δευτερόλεπτα--λεπτά & byte ανά συμβάν & 4G, PLC \\
Καμπύλες φορτίου μετρητών & ώρες & kB ανά μετρητή και ημέρα & 4G, PLC \\
\hline
\end{tabular}
\caption{Ενδεικτικές απαιτήσεις των εφαρμογών ενός διαχειριστή διανομής από το
τηλεπικοινωνιακό δίκτυο.}
\label{tab:comm-req}
\end{table}

\subsection{SCADA, DMS και ADMS}
\label{sec:adms}

Το SCADA των δικτύων διανομής συλλέγει τις μετρήσεις και τις ενδείξεις των υποσταθμών ΥΤ/ΜΤ,
των τηλεχειριζόμενων διακοπτών και των σταθμών παραγωγής και εκτελεί τους τηλεχειρισμούς,
όπως το SCADA του συστήματος μεταφοράς. Το \emph{σύστημα διαχείρισης δικτύου διανομής} (Distribution
Management System, DMS) είναι το σύνολο των εφαρμογών που επεξεργάζονται τα δεδομένα αυτά για
την ανάλυση και τη λειτουργία του δικτύου, όπως το EMS για το σύστημα μεταφοράς (Κεφάλαιο 1,
Ενότητα «\nameref{sec:meas-chain}»). Ιστορικά οι διαχειριστές διανομής λειτουργούσαν χωριστά
συστήματα: το SCADA, το σύστημα διαχείρισης διακοπών (Outage Management System, OMS), που
παρακολουθεί τις κλήσεις των καταναλωτών, εντοπίζει τις βλάβες και κατευθύνει τα συνεργεία,
και το σύστημα γεωγραφικών πληροφοριών (Geographic Information System, GIS), που περιέχει το
μοντέλο του δικτύου όπως κατασκευάστηκε. Το \emph{προηγμένο σύστημα διαχείρισης δικτύου διανομής}
(Advanced Distribution Management System, ADMS) ενοποιεί τα συστήματα αυτά σε μία πλατφόρμα με
κοινό μοντέλο δικτύου, το οποίο προέρχεται από το GIS, και κοινή διεπαφή χειριστή
(Σχήμα~\ref{fig:adms}), ώστε κάθε εφαρμογή να χρησιμοποιεί την ίδια τρέχουσα εικόνα του
δικτύου, και ανταλλάσσει δεδομένα με το MDMS, με το \emph{σύστημα διαχείρισης των ΔΕΠ}
(Distributed Energy Resources Management System, DERMS) και με τον ΔΣΜΗΕ. Οι κύριες
λειτουργίες του είναι οι εξής:

\begin{itemize}
\item \textbf{Διαχείριση του μοντέλου δικτύου.} Το μοντέλο εξάγεται από το GIS, όπου
      καταγράφονται τα στοιχεία του δικτύου, η συνδεσμολογία τους και οι καταναλωτές κάθε
      υποσταθμού ΜΤ/ΧΤ, και ενημερώνεται με τις θέσεις των διακοπτών από το SCADA και με τους
      χειροκίνητους χειρισμούς που καταχωρούν οι χειριστές. Η ανταλλαγή των μοντέλων μεταξύ
      των συστημάτων ακολουθεί το \emph{κοινό μοντέλο πληροφοριών} (Common Information Model,
      CIM) των προτύπων IEC 61970 και IEC 61968, το οποίο χρησιμοποιείται και για την
      ανταλλαγή δεδομένων με τον ΔΣΜΗΕ (Κεφάλαιο 7). Ο επεξεργαστής τοπολογίας προσδιορίζει σε κάθε
      στιγμή ποια τμήματα τροφοδοτούνται και από ποια γραμμή.
\item \textbf{Εκτίμηση κατάστασης (DSSE) και ροή ισχύος.} Υπολογίζουν τις τάσεις και τις ροές σε όλο
      το δίκτυο από τις λίγες μετρήσεις και τα προφίλ των φορτίων (Ενότητα
      «\nameref{sec:dsse}»), και παρέχουν την εικόνα του δικτύου στις υπόλοιπες εφαρμογές.
\item \textbf{Βελτιστοποίηση τάσης και άεργης ισχύος} (Volt/VAR Optimization, VVO). Καθορίζει
      τις τάσεις αναφοράς των ΣΑΤΥΦ, τη ζεύξη των πυκνωτών και την άεργη ισχύ των μονάδων
      παραγωγής, ώστε οι τάσεις να παραμένουν εντός ορίων με ελάχιστες απώλειες ή με
      μείωση της τάσης για εξοικονόμηση ενέργειας (Κεφάλαιο 2, Ενότητα «\nameref{sec:oltc}»).
\item \textbf{Εντοπισμός, απομόνωση σφάλματος και αποκατάσταση της τροφοδότησης} (Fault
      Location, Isolation and Service Restoration, FLISR), η οποία αναλύεται στην Ενότητα
      «\nameref{sec:flisr}».
\item \textbf{Διαχείριση διακοπών} (OMS). Συνδυάζει τις κλήσεις των καταναλωτών, τα μηνύματα των
      έξυπνων μετρητών και τις ενδείξεις των προστασιών για τον εντοπισμό της βλάβης,
      καταγράφει τους καταναλωτές που επηρεάζονται, κατευθύνει τα συνεργεία και υπολογίζει
      τους δείκτες αξιοπιστίας της \eqref{eq:saidi}.
\item \textbf{Διαχείριση χειρισμών και αναδιάταξη.} Καταρτίζει και ελέγχει τις ακολουθίες
      χειρισμών για τις προγραμματισμένες εργασίες και την αναδιάταξη, επαληθεύοντας με ροή
      ισχύος ότι δεν προκύπτουν υπερφορτίσεις ή αποκλίσεις τάσης.
\item \textbf{Πρόβλεψη φορτίου και παραγωγής} ανά γραμμή και υποσταθμό, σε ορίζοντα ωρών έως
      ημερών, από ιστορικά δεδομένα, δεδομένα των έξυπνων μετρητών και μετεωρολογικές
      προβλέψεις, ως είσοδος για τη διαχείριση των συμφορήσεων
      (Κεφάλαιο 6).
\end{itemize}

\begin{figure}[!htbp]
\centering
\includegraphics[width=\linewidth]{figures/adms.svg}
\caption{Αρχιτεκτονική της εποπτείας και του ελέγχου ενός δικτύου διανομής.}
\label{fig:adms}
\end{figure}

Η \emph{αναδιάταξη} για την ελαχιστοποίηση των απωλειών διατυπώνεται ως πρόβλημα
βελτιστοποίησης με μια δυαδική μεταβλητή $y_n \in \{0, 1\}$ για την κατάσταση κάθε κλάδου με
διακόπτη, κλειστού ή ανοικτού, το οποίο ελαχιστοποιεί τις απώλειες με περιορισμούς το πλήθος
των κλειστών κλάδων και τα όρια τάσης και ρεύματος:

\begin{equation}
\begin{aligned}
& \min_{\mathbf{y}} \sum_{n} y_n\, r_n\,\frac{P_n^2 + Q_n^2}{V_{\pi(n)}^2} \\
& \sum_{n} y_n = N_k - N_s, \qquad V^{\min} \le V_k \le V^{\max}, \qquad I_n \le I_n^{\max}
\end{aligned}
\label{eq:reconf}
\end{equation}

όπου οι ροές $P_n$, $Q_n$ και οι τάσεις προκύπτουν από τις εξισώσεις DistFlow του Κεφαλαίου 5
για τους κλειστούς κλάδους, $N_k$ το πλήθος των κόμβων και $N_s$ το πλήθος των σημείων
τροφοδότησης, δηλαδή των ζυγών ΜΤ. Ο περιορισμός του πλήθους των κλειστών κλάδων, μαζί με την
απαίτηση να είναι το δίκτυο συνεκτικό, εξασφαλίζει ότι κάθε κόμβος τροφοδοτείται από μία
ακριβώς διαδρομή, δηλαδή την ακτινική λειτουργία.

Το DERMS συμπληρώνει το ADMS με τις λειτουργίες που αφορούν τους ΔΕΠ: καταγραφή και
παρακολούθηση των μονάδων, πρόβλεψη της παραγωγής και της ευελιξίας τους, ομαδοποίησή τους ανά
γραμμή ή υποσταθμό, αποστολή εντολών, όπως ορίων ενεργής ισχύος ή
τιμών αναφοράς άεργης ισχύος, είτε απευθείας είτε μέσω φορέων σωρευτικής εκπροσώπησης
(aggregators), και επαλήθευση της απόκρισης. Ο οδηγός IEEE 2030.11 ορίζει τις λειτουργικές
προδιαγραφές ενός DERMS, ενώ η επικοινωνία με τις μονάδες χρησιμοποιεί πρωτόκολλα όπως το
IEC 61850 και το IEEE 2030.5. Το DERMS είναι ο δίαυλος μέσω του οποίου ο διαχειριστής διανομής
αξιοποιεί την ευελιξία των ΔΕΠ για τη διαχείριση των συμφορήσεων
(Κεφάλαιο 6) και τη διαθέτει στον ΔΣΜΗΕ
(Κεφάλαιο 7).

\subsection{Εκτίμηση κατάστασης στα δίκτυα διανομής}
\label{sec:dsse}

Η \emph{εκτίμηση κατάστασης δικτύων διανομής} (Distribution System State Estimation, DSSE)
προσδιορίζει τις τάσεις όλων των κόμβων του δικτύου από τις διαθέσιμες μετρήσεις, όπως η
εκτίμηση κατάστασης του συστήματος μεταφοράς (Κεφάλαιο 1, Ενότητα «\nameref{sec:se}»), και
είναι η βάση κάθε εφαρμογής του ADMS που χρειάζεται την τρέχουσα εικόνα του δικτύου: του
VVO, της διαχείρισης των συμφορήσεων, του FLISR και του υπολογισμού της ικανότητας υποδοχής. Η διατύπωση είναι η ίδια: με
διάνυσμα μετρήσεων $\mathbf{z}$, διάνυσμα κατάστασης $\mathbf{x}$ και συναρτήσεις μετρήσεων
$\mathbf{h}(\mathbf{x})$, η εκτίμηση σταθμισμένων ελαχίστων τετραγώνων ελαχιστοποιεί το

\begin{equation}
J(\mathbf{x}) = \sum_{i=1}^{m} \frac{\left(z_i - h_i(\mathbf{x})\right)^2}{\sigma_i^2}
\label{eq:dsse-wls}
\end{equation}

όπου $\sigma_i^2$ η διακύμανση του σφάλματος της μέτρησης $i$, με επαναληπτική επίλυση των
κανονικών εξισώσεων όπως στην Ενότητα «\nameref{sec:wls}» του Κεφαλαίου 1. Οι διαφορές
βρίσκονται στα δεδομένα και στο μοντέλο (Πίνακας~\ref{tab:dsse}).

\begin{table}[!htbp]
\centering
\begin{tabular}{lll}
\hline
 & Σύστημα μεταφοράς & Δίκτυο διανομής \\
\hline
Μετρήσεις πραγματικού χρόνου & σε όλους τους ζυγούς και κλάδους & υποσταθμός ΥΤ/ΜΤ, λίγα σημεία ΜΤ \\
Πλεονασμός μετρήσεων & 2--3 & $<$ 1 χωρίς ψευδομετρήσεις \\
Ψευδομετρήσεις & σπάνια & σε κάθε κόμβο φορτίου \\
Συγχρονισμός & SCADA (s), PMU (ms) & SCADA· έξυπνοι μετρητές ανά 15 min \\
Μοντέλο δικτύου & μονοφασικό ισοδύναμο & τριφασικό, ασύμμετρο \\
Μεταβλητές κατάστασης & τάσεις ζυγών & τάσεις κόμβων ανά φάση \\
Πλήθος κόμβων & εκατοντάδες & χιλιάδες ανά υποσταθμό \\
Ανίχνευση εσφαλμένων δεδομένων & με τον πλεονασμό & περιορισμένη \\
\hline
\end{tabular}
\caption{Διαφορές της εκτίμησης κατάστασης στο σύστημα μεταφοράς και στα δίκτυα διανομής.}
\label{tab:dsse}
\end{table}

Η κύρια διαφορά είναι ο \emph{πλεονασμός} των μετρήσεων. Στο σύστημα μεταφοράς οι μετρήσεις
είναι δύο έως τρεις φορές περισσότερες από τις μεταβλητές κατάστασης, οπότε η εκτίμηση
εξομαλύνει τα σφάλματα και ανιχνεύει τα εσφαλμένα δεδομένα. Στο δίκτυο διανομής οι
μετρήσεις πραγματικού χρόνου περιορίζονται στην αναχώρηση κάθε γραμμής στον υποσταθμό ΥΤ/ΜΤ,
δηλαδή στην τάση των ζυγών ΜΤ και στο ρεύμα ή στην ενεργή και άεργη ισχύ της γραμμής, και σε
λίγα σημεία κατά μήκος της, όπως οι τηλεχειριζόμενοι διακόπτες και οι διακόπτες αυτόματης
επαναφοράς, που μετρούν συνήθως μόνο ρεύμα, καθώς και στους μεγάλους σταθμούς παραγωγής. Το
δίκτυο δεν είναι παρατηρήσιμο με αυτές μόνο τις μετρήσεις (Κεφάλαιο 1, Ενότητα
«\nameref{sec:observability}»). Η παρατηρησιμότητα αποκαθίσταται με \emph{ψευδομετρήσεις}
(pseudo-measurements): εκτιμήσεις της ενεργής και άεργης ισχύος κάθε κόμβου φορτίου από
στατιστικά προφίλ φορτίου ανά κατηγορία καταναλωτή, ώρα και εποχή, κλιμακωμένα με την
ετήσια κατανάλωση που είναι γνωστή από την τιμολόγηση, ή από πρόβλεψη φορτίου. Οι
ψευδομετρήσεις εισάγονται στην \eqref{eq:dsse-wls} με μεγάλη διακύμανση, που αντιστοιχεί σε
σφάλμα 20--50\%, ώστε να μην αλλοιώνουν τις πραγματικές μετρήσεις, οι οποίες έχουν σφάλμα
1--3\%. Οι \emph{εικονικές μετρήσεις} των κόμβων χωρίς φορτίο, όπως οι κόμβοι διακλάδωσης,
είναι μηδενικές εγχύσεις που είναι γνωστές με ακρίβεια και εισάγονται με πολύ μικρή
διακύμανση. Με τις ψευδομετρήσεις η εκτίμηση συγκλίνει πάντοτε, αλλά η ακρίβειά της στους
κόμβους μακριά από τις πραγματικές μετρήσεις καθορίζεται από την ακρίβεια των προφίλ. Το
Σχήμα~\ref{fig:dsse} συνοψίζει τα είδη των δεδομένων σε μια γραμμή ΜΤ: τις
μετρήσεις πραγματικού χρόνου στην αναχώρηση της γραμμής και στους τηλεχειριζόμενους διακόπτες,
τις μετρήσεις των σταθμών παραγωγής που εντάσσονται στο SCADA, τις ψευδομετρήσεις των φορτίων
των υποσταθμών ΜΤ/ΧΤ, τις εικονικές μετρήσεις στους κόμβους διακλάδωσης και τα αθροίσματα των
έξυπνων μετρητών ανά υποσταθμό ΜΤ/ΧΤ, τα οποία είναι διαθέσιμα ανά 15 λεπτά.

\begin{figure}[!htbp]
\centering
\includegraphics[width=\linewidth]{figures/dsse-measurements.svg}
\caption{Είδη δεδομένων της εκτίμησης κατάστασης σε γραμμή ΜΤ.}
\label{fig:dsse}
\end{figure}

Οι έξυπνοι μετρητές βελτιώνουν ουσιαστικά την εκτίμηση. Τα δεδομένα τους δεν είναι όμως
συγχρονισμένα με το SCADA ούτε διαθέσιμα σε πραγματικό χρόνο: αναφέρονται σε μέσες
τιμές διαστήματος και φθάνουν στον διαχειριστή με καθυστέρηση λεπτών έως ωρών, οπότε
συνδυάζονται με τις μετρήσεις πραγματικού χρόνου είτε ως ψευδομετρήσεις βελτιωμένης
ακρίβειας, με κατάλληλη στάθμιση, είτε για την επικαιροποίηση των προφίλ φορτίου.

Η δεύτερη διαφορά είναι το \emph{μοντέλο}. Τα δίκτυα ΧΤ είναι ασύμμετρα και απαιτούν
\emph{τριφασική εκτίμηση κατάστασης}. Ένα τμήμα γραμμής
τεσσάρων αγωγών, με τις φάσεις a, b, c και τον ουδέτερο n, περιγράφεται από τον πίνακα $4
\times 4$ των ιδίων και των αμοιβαίων σύνθετων αντιστάσεών του. Αν ο ουδέτερος είναι γειωμένος
στα δύο άκρα του τμήματος, όπως στην ουδετέρωση (σύστημα TN) της ΧΤ, η τάση του ως προς γη
θεωρείται μηδενική και η εξίσωσή του απαλείφεται με τη \emph{μείωση Kron}:

\begin{equation}
\mathbf{Z}_{abc} = \mathbf{Z}_{pp} - \mathbf{Z}_{pn}\,Z_{nn}^{-1}\,\mathbf{Z}_{np}
\label{eq:kron}
\end{equation}

όπου $\mathbf{Z}_{pp}$ ο πίνακας $3 \times 3$ των φάσεων, $\mathbf{Z}_{pn} = \mathbf{Z}_{np}^{T}$
η στήλη των αμοιβαίων σύνθετων αντιστάσεων φάσεων-ουδετέρου και $Z_{nn}$ η ιδία σύνθετη
αντίσταση του ουδετέρου. Ο \emph{πίνακας φάσεων} $\mathbf{Z}_{abc}$ συνδέει τις πτώσεις τάσης με
τα ρεύματα των τριών φάσεων, $\Delta\bar{\mathbf{V}}_{abc} = \mathbf{Z}_{abc}\,\bar{\mathbf{I}}_{abc}$,
και τα μη διαγώνια στοιχεία του εκφράζουν την επίδραση του ρεύματος κάθε φάσης, μέσω της
επιστροφής από τον ουδέτερο και της αμοιβαίας επαγωγής, στην τάση των άλλων. Για γραμμή με
συμμετρική διάταξη των αγωγών, όπως τα καλώδια και οι γραμμές με μεταθέσεις, ο πίνακας
προκύπτει από τις σύνθετες αντιστάσεις ορθής ακολουθίας $z_1$ και ομοπολικής ακολουθίας
$z_0$:

\begin{equation}
\mathbf{Z}_{abc} =
\begin{bmatrix} z_s & z_m & z_m \\ z_m & z_s & z_m \\ z_m & z_m & z_s \end{bmatrix}, \qquad
z_s = \frac{z_0 + 2z_1}{3}, \qquad z_m = \frac{z_0 - z_1}{3}
\label{eq:zabc-sym}
\end{equation}

Με τον πίνακα αγωγιμοτήτων κόμβων, του οποίου κάθε στοιχείο είναι πλέον ένα μπλοκ
$3 \times 3$ που σχηματίζεται από τους αντίστροφους πίνακες $\mathbf{Z}_{abc}^{-1}$ των
τμημάτων, η έγχυση ισχύος στη φάση $\phi$ του κόμβου $k$ είναι

\begin{equation}
P_k^{\phi} + \mathrm{j}Q_k^{\phi}
  = \bar{V}_k^{\phi} \sum_{j} \sum_{\psi \in \{a, b, c\}}
    \left(\bar{Y}_{kj}^{\phi\psi}\,\bar{V}_j^{\psi}\right)^{*}
\label{eq:dsse-3f}
\end{equation}

όπου $\bar{Y}_{kj}^{\phi\psi}$ το στοιχείο της γραμμής $\phi$ και της στήλης $\psi$ του μπλοκ
$(k, j)$. Οι συναρτήσεις μετρήσεων των ροών και των ρευμάτων ανά φάση σχηματίζονται με τον ίδιο
τρόπο, και κάθε μέτρηση έξυπνου μετρητή αντιστοιχεί σε συγκεκριμένο κόμβο και φάση. Οι
μεταβλητές κατάστασης είναι το μέτρο και η γωνία της τάσης κάθε φάσης σε κάθε κόμβο, δηλαδή
$6N_k - 3$ για $N_k$ κόμβους: οι τρεις γωνίες του υποσταθμού λαμβάνονται ως αναφορά, 0, −120°
και 120°.

\subsection{Αυτόματη αποκατάσταση της τροφοδότησης}
\label{sec:flisr}

Η λειτουργία \emph{εντοπισμού, απομόνωσης σφάλματος και αποκατάστασης της τροφοδότησης}
(FLISR) αυτοματοποιεί τα δύο πρώτα στάδια της αντιμετώπισης μιας διακοπής και είναι η
εφαρμογή του ADMS με το πιο άμεσο όφελος στους δείκτες αξιοπιστίας. Προϋποθέτει γραμμές ΜΤ
χωρισμένες σε τμήματα με τηλεχειριζόμενους διακόπτες φορτίου ή διακόπτες αυτόματης
επαναφοράς, \emph{ενδεικτικά διέλευσης σφάλματος} (Fault Passage Indicators, FPI) με
τηλεένδειξη, τα οποία αναφέρουν στο κέντρο ελέγχου ότι το ρεύμα σφάλματος πέρασε από το σημείο
εγκατάστασής τους, και ΚΑ προς γειτονικές γραμμές. Μετά από μόνιμο σφάλμα σε ένα τμήμα, ο
διακόπτης ισχύος της γραμμής ανοίγει και όλοι οι καταναλωτές της διακόπτονται. Η λειτουργία εκτελεί
στη συνέχεια τα εξής βήματα:

\begin{enumerate}
\item \textbf{Εντοπισμός.} Από τις ενδείξεις των FPI και των προστασιών, το σφάλμα
      εντοπίζεται στο τμήμα μεταξύ του τελευταίου FPI που ενεργοποιήθηκε και του πρώτου που δεν ενεργοποιήθηκε, ενώ η θέση του μέσα στο τμήμα
      εκτιμάται από τις τάσεις και τα ρεύματα που κατέγραψε ο ηλεκτρονόμος της γραμμής.
\item \textbf{Απομόνωση.} Οι τηλεχειριζόμενοι διακόπτες στα δύο άκρα του τμήματος που έπαθε
      βλάβη ανοίγουν.
\item \textbf{Αποκατάσταση ανάντη.} Ο διακόπτης ισχύος της γραμμής κλείνει και
      επανατροφοδοτεί τους καταναλωτές μεταξύ του υποσταθμού και του σφάλματος.
\item \textbf{Αποκατάσταση κατάντη.} Οι καταναλωτές κατάντη του σφάλματος επανατροφοδοτούνται
      από γειτονική γραμμή με το κλείσιμο του ΚΑ, αφού η ροή ισχύος
      επαληθεύσει ότι η γειτονική γραμμή και ο μετασχηματιστής του υποσταθμού της αντέχουν
      το πρόσθετο φορτίο και ότι οι τάσεις παραμένουν εντός ορίων. Αν το φορτίο υπερβαίνει τη
      διαθέσιμη ικανότητα, αποκαθίσταται μέρος των καταναλωτών ή χρησιμοποιούνται
      περισσότερες γειτονικές γραμμές.
\end{enumerate}

Στο Σχήμα~\ref{fig:flisr}, μετά από σφάλμα στο τρίτο τμήμα της γραμμής Α, ο διακόπτης ισχύος
της γραμμής ανοίγει (1), οι τηλεχειριζόμενοι διακόπτες στα άκρα του τμήματος απομονώνουν το
σφάλμα (2), ο διακόπτης της γραμμής κλείνει και επανατροφοδοτεί τα τμήματα ανάντη (3) και ο
ΚΑ προς τη γραμμή Β κλείνει και επανατροφοδοτεί το τμήμα κατάντη (4).
Χωρίς τροφοδότηση παραμένει μόνο το τμήμα με το σφάλμα, έως την επισκευή του.

\begin{figure}[!htbp]
\centering
\includegraphics[width=\linewidth]{figures/flisr.svg}
\caption{Στάδια της αυτόματης αποκατάστασης της τροφοδότησης.}
\label{fig:flisr}
\end{figure}

Η ακολουθία εκτελείται είτε κεντρικά από το ADMS, το οποίο διαθέτει το πλήρες μοντέλο του
δικτύου και μπορεί να ελέγξει τις ροές ισχύος και τις τάσεις πριν από κάθε χειρισμό, είτε
αποκεντρωμένα από τις IED των διακοπτών, που επικοινωνούν μεταξύ τους με μηνύματα GOOSE και
αποκρίνονται ταχύτερα αλλά με απλούστερους κανόνες. Σε κάθε περίπτωση οι καταναλωτές των
τμημάτων χωρίς βλάβη επανατροφοδοτούνται εντός ενός έως λίγων λεπτών, αντί για τη μισή έως
δύο ώρες που απαιτούν οι χειροκίνητοι χειρισμοί, οπότε η διακοπή τους δεν προσμετράται καν
στις διακοπές μακράς διάρκειας του SAIDI. Το όφελος εξαρτάται επομένως από το πλήθος
των τμημάτων κάθε γραμμής και από τη διαθέσιμη ικανότητα των γειτονικών γραμμών· το δεύτερο
παράδειγμα του κεφαλαίου υπολογίζει τους δείκτες αξιοπιστίας μιας γραμμής για
διαφορετικούς βαθμούς αυτοματισμού. Η διανεμημένη παραγωγή επηρεάζει τη λειτουργία με δύο
τρόπους: οι σταθμοί κατάντη του σφάλματος πρέπει να αποσυνδέονται πριν από την απομόνωση και
την επανατροφοδότηση, ώστε να αποφεύγεται η ακούσια νησιδοποίηση και η επανασύνδεση εκτός
συγχρονισμού (Κεφάλαιο 2, Ενότητα «\nameref{sec:island}»), ενώ η συμβολή τους στο ρεύμα
σφάλματος μπορεί να ενεργοποιήσει FPI σε τμήματα χωρίς βλάβη, οπότε τα FPI πρέπει να είναι
κατευθυντικά.
